\setcounter{aufgabe}{21}
\einheitenkopf{Einheit 3 von 5 · Streifen- und Kreisdiagramm}

\begin{aufgabe}{Anteile -- Prozent oder Grad? Kreuze an, welche Einheit in die Lücke gehört. Rechne nicht.}
\begin{teile}
\teil Ein Viertel des Kreises ist ein Sektor von 90 \underline{\hspace{1cm}} \hfill \kreuz{\%}\kreuz{$^\circ$}
\teil Ein Viertel aller Befragten sind 25 \underline{\hspace{1cm}} \hfill \kreuz{\%}\kreuz{$^\circ$}
\teil Der ganze Kreis misst 360 \underline{\hspace{1cm}} \hfill \kreuz{\%}\kreuz{$^\circ$}
\teil Alle Anteile eines Kreisdiagramms ergeben zusammen 100 \underline{\hspace{1cm}} \hfill \kreuz{\%}\kreuz{$^\circ$}
\teil Von 200 Personen nennen 40 den Bus; ihr Sektor misst 72 \underline{\hspace{1cm}} \hfill \kreuz{\%}\kreuz{$^\circ$}
\end{teile}
\end{aufgabe}

\begin{aufgabe}{Anteile -- am Streifen. Der ganze Streifen ist 10\,cm lang und steht für $100\,\%$. Wie lang ist der Abschnitt?}
\beispiel{30\,\% \quad &\rightarrow \quad 3\,\text{cm}}
\begin{teilezwei}
\tz $40\,\%$ \leerfeld[cm] & \tz $20\,\%$ \leerfeld[cm] \\
\tz $70\,\%$ \leerfeld[cm] & \tz $5\,\%$ \leerfeld[cm] \\
\tz $45\,\%$ \leerfeld[cm] & \\
\end{teilezwei}
\begin{teile}
\teil Trage im Streifen die vier Abschnitte ab und beschrifte sie: Rad $25\,\%$, Bus $40\,\%$, zu Fuß $20\,\%$, Auto $15\,\%$.
\end{teile}

\noindent\streifenleer

\begin{teile}
\teil In einem anderen Streifen sind bereits $30\,\%$, $25\,\%$ und $20\,\%$ eingetragen. Wie viel Prozent fehlen noch? \leerfeld[\%]
\end{teile}
\end{aufgabe}

\begin{aufgabe}{Anteile -- Winkel berechnen. Berechne den Mittelpunktswinkel des Sektors.}
\beispiel{30\,\% \quad &\rightarrow \quad 0{,}30 \cdot 360^\circ = 108^\circ}
\begin{teilezwei}
\tz $50\,\%$ \leerfeld[$^\circ$] & \tz $25\,\%$ \leerfeld[$^\circ$] \\
\tz $10\,\%$ \leerfeld[$^\circ$] & \tz $75\,\%$ \leerfeld[$^\circ$] \\
\end{teilezwei}
\begin{teile}
\teil $12\,\%$ \leerfeld[$^\circ$]
\teil $37{,}5\,\%$ \leerfeld[$^\circ$]
\teil 8 von 25 Personen \leerfeld[$^\circ$]
\teil 15 von 400 Personen \leerfeld[$^\circ$]
\teil 3 von 7 Personen, gerundet auf eine Stelle nach dem Komma \leerfeld[$^\circ$]
\end{teile}
\end{aufgabe}

\begin{aufgabe}{Anteile -- Kreisdiagramm zeichnen. Nutze für a) bis c) den Kreis mit dem eingezeichneten Radius.}
\begin{teile}
\teil Trage vom Radius aus mit dem Geodreieck einen Sektor von $90^\circ$ an.
\teil Trage daran anschließend einen Sektor von $120^\circ$ an.
\teil Wie groß ist der Winkel des Restsektors? \leerfeld[$^\circ$]
\end{teile}

\noindent\kreisleer[2.2]

\begin{teile}
\teil Ein Verein hat 240 Mitglieder: 100 spielen Fußball, 80 Handball, 60 machen Judo. Berechne die drei Mittelpunktswinkel.
\par Fußball \leerfeld[$^\circ$] \quad Handball \leerfeld[$^\circ$] \quad Judo \leerfeld[$^\circ$]
\teil Zeichne damit das Kreisdiagramm und beschrifte die drei Sektoren.
\end{teile}

\noindent\kreisleer[2.2]
\end{aufgabe}

\begin{aufgabe}{Anteile -- Sektoren zuordnen. Das Kreisdiagramm zeigt, womit 400 Haushalte einer Stadt heizen: Gas $38\,\%$, Öl $27\,\%$, Fernwärme $21\,\%$, Wärmepumpe $14\,\%$. Die Sektoren sind nicht beschriftet.}
\noindent\kreisdiagrammleer[2]{/38, /27, /21, /14}

\begin{teile}
\teil Ordne die vier Angaben den Sektoren zu und schreibe sie auf die Striche.
\teil Welcher Sektor ist ungefähr ein Viertel des Kreises? \leerfeld
\teil In einer zweiten Stadt heizen 17 von 450 Haushalten mit Holz. Berechne den Mittelpunktswinkel dieses Sektors und runde auf eine Stelle nach dem Komma. (P10 2024) \leerfeld[$^\circ$]
\end{teile}
\end{aufgabe}

\begin{aufgabe}{Anteile -- Fehler finden.}
Für ein Kreisdiagramm soll der Sektor für $24\,\%$ gezeichnet werden. Ben schreibt auf:
\rechnung{\text{Sektor}\ &= 24^\circ}
\begin{teile}
\teil Finde den Fehler, benenne ihn und berechne den richtigen Winkel. \\[12pt]
\teil Berechne den Mittelpunktswinkel für einen Anteil von $45\,\%$. \leerfeld[$^\circ$]
\end{teile}
\end{aufgabe}

\begin{aufgabe}{Anteile -- Begründen.}
\begin{teile}
\teil Erkläre, warum ein Anteil von $1\,\%$ im Kreisdiagramm genau $3{,}6^\circ$ entspricht. \\[12pt]
\teil In einem Kreisdiagramm ist ein Sektor größer als der halbe Kreis. Was weißt du über seinen Anteil, ohne zu messen? Begründe. \\[12pt]
\end{teile}
\end{aufgabe}
